Sources : DGFiP, déclarations 2042 des revenus de 2024 et Statistiques n° 53 (revenus de 2025) ; Lefebvre, Lehmann et Sicsic (2025) et CAE (2025) pour l'élasticité ; calculs de l'auteur (pfu/calculs\_pfu.py).\par Note : revenus de 2025 (assiette estimée à 67 Md\euro{}). Traits : élasticité des revenus du capital au taux net de 0,65 à 1 (0,77 au centre). Les chiffres portent sur les premières années, où toute la réaction compte comme recette. En régime permanent, si la réaction n'est qu'un décalage dans le temps, un PFU à 25 \% coûte jusqu'à 4,3 Md\euro{} et un PFU à 33 \% rapporte jusqu'à 1,1 Md\euro{}. Sur les revenus de 2024, gonflés par l'anticipation de la contribution différentielle, ces montants seraient de 5,1 et 1,3 Md\euro{}. Retour au barème : ordre de grandeur sous hypothèses fortes (annexe).
